\section*{Capítulo \ref{ch:g4:what-plants-need} --- Lo que necesitan las plantas}

\begin{solution}{exo:g4:what-plants-need:1}
Agua, luz, aire, calor y \omterm{prop:g4:what-plants-need:needs}{sales minerales} (bebidas disueltas en el agua
de la tierra).
\end{solution}

\begin{solution}{exo:g4:what-plants-need:2}
Exactamente una --- la que se está probando; todo lo demás tiene que
ser idéntico. La maceta que no se toca es el testigo.
\end{solution}

\begin{solution}{exo:g4:what-plants-need:3}
Amarilla pálida en vez de verde, raramente alta y estirada, débil ---
y después se rinde. Faltaba la luz, así que la \omterm{def:g1:plants-around-us:plant}{planta} no pudo fabricar
su sustancia verde ni su alimento.
\end{solution}

\begin{solution}{exo:g4:what-plants-need:4}
Cantidades diminutas de sustancias nutritivas que hay en la tierra.
Las \omterm{def:g1:plants-around-us:parts}{raíces} las beben disueltas en el agua de la tierra.
\end{solution}

\begin{solution}{exo:g4:what-plants-need:5}
Porque las cosechas se llevan las \omterm{prop:g4:what-plants-need:needs}{sales minerales} de la tierra año
tras año; el compost las repone. Alimenta a la \emph{tierra}, y el
agua de la tierra lleva después las sales hasta las \omterm{def:g1:plants-around-us:parts}{raíces}.
\end{solution}

\begin{solution}{exo:g4:what-plants-need:6}
Un \omterm{def:g1:animals-around-us:animal}{animal} crece a partir de los \omterm{def:g1:living-or-not:living}{seres vivos} que se come; una \omterm{def:g1:plants-around-us:plant}{planta} no
se come a nadie --- fabrica su propia materia con agua, \omterm{prop:g4:what-plants-need:needs}{sales
minerales}, aire y la \omterm{prop:g3:balanced-diet:jobs}{energía} de la luz que atrapan sus \omterm{def:g1:plants-around-us:parts}{hojas} verdes.
\end{solution}

\begin{solution}{exo:g4:what-plants-need:7}
De la tierra no: del agua, del aire, y con la luz aportando la
\omterm{prop:g3:balanced-diet:jobs}{energía} --- construida en las \omterm{def:g1:plants-around-us:parts}{hojas} verdes. (La tierra solo puso agua
y pequeñas cantidades de \omterm{prop:g4:what-plants-need:needs}{sales minerales}.)
\end{solution}

\begin{solution}{exo:g4:what-plants-need:8}
Respuesta abierta. Lo esperable: la maceta regada verde y erguida, la
gemela seca marchitándose en unos días --- y como lo único distinto
era el agua, el mérito se lo lleva el agua.
\end{solution}

\begin{solution}{exo:g4:what-plants-need:9}
No hay testigo: una sola maceta y todo a la vez --- puede que
sencillamente tuviera buena luz y un riego cuidadoso (¡quien le canta
a una \omterm{def:g1:plants-around-us:plant}{planta} la visita muy a menudo!). La versión justa: dos macetas
idénticas, con las mismas \omterm{prop:g3:flowers-fruits-seeds:transform}{semillas}, la misma tierra, la misma ventana
y el mismo riego; cantarle solo a una; comparar al cabo de un mes.
Solo entonces una diferencia apuntaría al canto.
\end{solution}

\begin{solution}{exo:g4:what-plants-need:10}
La luz: crece sumergida, pero en agua iluminada, y sus \omterm{def:g1:plants-around-us:parts}{hojas} verdes
siguen atrapando la luz a través del agua. El agua que la rodea le
aporta además lo que aporta el agua de la tierra --- \omterm{prop:g4:what-plants-need:needs}{sales minerales}
disueltas, bebidas por toda su superficie en vez de por unas \omterm{def:g1:plants-around-us:parts}{raíces}.
\end{solution}

\begin{solution}{exo:g4:what-plants-need:11}
Las \omterm{prop:g4:what-plants-need:needs}{sales minerales}: el agua, la luz, el aire y el calor eran iguales,
y la arena pura con agua de lluvia casi no aporta sales. Al principio
creció con las reservas guardadas en sus propios tejidos (como la
\omterm{prop:g2:from-seed-to-plant:cycle}{plántula} enterrada crecía con la despensa de la \omterm{def:g2:from-seed-to-plant:seed}{semilla}) --- y
gastadas las reservas, se notó la falta.
\end{solution}
